% =============================================================================
% Project : THE_PHILOSOPHY_OF_VIABILITY
% Book    : Book I — Ontology of Viability
% Book ID : VS.B01
%
% File    : 2_MAINMATTER/PART_I_IDENTITY_AND_MANIFESTATION/CHAPTER_02_INVARIANCE/sections/section_02_invariance_and_identity.tex
%
% Section / Chapter:
%   Section 02 — Invariance and Identity
%   Chapter 02 — Invariance
%
% Purpose :
%   To establish the ontological distinction between identity and
%   invariance, examine the limits of invariant preservation as a
%   criterion of identity, and clarify the possible relevance of
%   invariance to the characterization of an existent.
%
% Conceptual Policy:
%   1. Being and identity are ontologically inseparable but
%      analytically distinguishable.
%   2. Neither being nor identity has ontological priority over
%      the other.
%   3. Invariance is neither the definition nor the foundation
%      of identity.
%   4. The preservation of an invariant does not, by itself,
%      establish numerical identity.
%   5. The absence of an identified invariant does not establish
%      the absence of identity.
%   6. Multiple invariants do not entail multiple identities,
%      nor do they establish a single identity.
%   7. The possible plurality and relations of identities remain
%      open ontological questions.
%
% =============================================================================

\section{Invariance and Identity}
\label{sec:invariance-and-identity}

\subsection{The Need for an Ontological Distinction}

Invariance and identity become relevant to the same
ontological problem whenever an existent is considered
in relation to transformation. The analysis of
transformation may raise two questions: what, if
anything, remains preserved, and whether the existent
in question retains its identity.

These questions are related, but they are not identical.

Invariance concerns the preservation of a specified
determination, relation, or configuration under a
specified transformation and within a specified domain.
Identity concerns the ontological identity of the
existent. The former examines a possible relation
between a feature and a transformation; the latter
concerns the existent in its identity.

The distinction is not merely terminological. Each
question requires a different kind of ontological
clarification. Establishing that a determination
remains preserved does not, without further argument,
establish what the identity of the existent consists
in. Conversely, establishing that a determination
has changed does not, by itself, settle whether the
existent retains its identity.

A preserved feature may be relevant to an account
of an existent, but relevance must not be confused
with constitution or sufficiency. The fact that a
determination remains unchanged does not establish
that it constitutes the identity of the existent,
and the fact that several determinations are preserved
does not establish that their collection is identical
with that identity.

The ontological distinction must therefore be maintained
before any account of the relation between identity
and invariance can be developed. Without it, the
analysis risks replacing the problem of identity with
a narrower inquiry into the preservation of selected
features.

\subsection{The Ontological Status of Identity}

The preceding chapter established three distinct
principles concerning the relation between being
and identity: their analytical distinction, their
ontological inseparability, and the absence of
ontological priority between them.

These principles must remain distinct. Analytical
distinction permits being and identity to be examined
through different questions. Ontological inseparability
denies that they constitute independently existing
realities. The absence of ontological priority denies
that either can be established as ontologically
prior to the other.

Within this framework, identity is not an additional
property attached to an existent after its existence
has already been constituted. Nor is it a secondary
result produced by the preservation of particular
determinations.

An existent is not first ontologically complete
and subsequently supplied with identity through
the maintenance of an invariant configuration.
Such an account would make identity dependent upon
a process of preservation and would thereby reduce
an ontological question to a condition concerning
the behaviour of selected features.

Identity must instead be considered in relation
to the existent itself, without presupposing that
its ontological status can be derived from a list
of preserved determinations.

This does not mean that identity has ontological
priority over being. The distinction between the
ontological status of an existent and the preservation
of its features does not introduce a hierarchy between
being and identity. It establishes only that invariance
and identity address different questions.

Invariance concerns whether a specified feature
is preserved under a transformation. It does not
constitute an independent ontological ground from
which the identity of the existent can be derived.

\subsection{The Insufficiency of Invariance as a Criterion of Identity}

To reduce identity to invariance would be to maintain
that the identity of an existent is determined by
the feature or features that remain preserved under
a relevant class of transformations.

Such a reduction requires an additional premise:
that the preserved features selected for analysis
are constitutive of, or sufficient to establish,
the identity of the existent.

That premise cannot be obtained from preservation
alone.

A feature may remain unchanged without being
constitutive of the identity of the existent. A
quantitative determination may remain constant
while the composition of an existent changes.
A structural relation may be preserved while
particular components are replaced. A qualitative
determination may remain unchanged while other
aspects of the existent undergo substantial
modification.

These cases do not establish that identity is
unaffected by every such change. They establish
a more limited point: preservation of a feature
does not, by itself, determine the ontological
significance of that feature for identity.

The selection of an invariant as an identity criterion
therefore requires justification beyond the fact
that the feature remains preserved. It must be
established why that particular determination
or relation should be regarded as relevant to
identity, and why its preservation should be
sufficient for the intended conclusion.

Without such justification, the reduction remains
unsupported.

A further difficulty arises from the plurality
of possible invariants. Different determinations
and relations may remain preserved under different
transformations, or within different domains of
the same transformation. A quantitative feature,
a structural relation, and a relational configuration
may each exhibit preservation under specified
conditions.

Their preservation does not establish that they
are equivalent in ontological significance. Nor
does it provide a rule for deciding whether one
of them, several of them, or their combination
should count as constitutive of identity.

The problem cannot be resolved merely by collecting
more invariants. An aggregate of preserved features
does not become an account of identity simply
because the number of preserved features increases.

The reduction of identity to invariance therefore
faces a problem of justification: preservation
establishes a fact about a specified feature under
specified conditions, but it does not independently
establish the ontological role of that feature
in the identity of an existent.

\subsection{The Preservation of an Invariant and Numerical Identity}

Numerical identity concerns whether two designations
or configurations refer to one and the same existent,
rather than to two distinct existents that may be
qualitatively similar, structurally corresponding,
or otherwise alike.

The preservation of an invariant does not, by itself,
resolve this question.

Two numerically distinct existents may share the
same quantitative determination, possess corresponding
structural relations, or exhibit equivalent patterns
of organization. A feature may therefore be preserved
or shared across distinct existents without making
them numerically identical.

The preservation of a feature establishes sameness
in a specified respect. Numerical identity is not
equivalent to sameness in every respect, nor can
it be inferred from the preservation of one selected
respect.

The same distinction applies when an existent is
considered across transformation. If a determination
remains preserved between an initial and a resulting
configuration, this establishes preservation with
respect to that determination. It does not, without
a further ontological argument, establish that
the configurations belong to one and the same
existent.

The relation between configurations and the identity
of an existent must therefore be examined independently
of the preservation claim.

This distinction does not make invariance irrelevant
to the investigation of numerical identity. An
invariant may contribute to an argument concerning
the relation between configurations. It may provide
a relevant condition, constraint, or indication
within a broader account. What it cannot do is
establish numerical identity merely by being
preserved.

The evidential or explanatory relevance of an
invariant must consequently be distinguished from
its sufficiency as a criterion of numerical identity.

\subsection{The Absence of an Identified Invariant and the Question of Identity}

The converse reduction is equally unwarranted.
If identity were defined through invariant
preservation, the failure to identify an invariant
would appear to undermine the identity of the
existent.

That conclusion does not follow.

The absence of an identified invariant may mean
that no relevant feature has been shown to remain
preserved within the specified domain and under
the transformation examined. It may also indicate
that the selected respect is not preserved, that
the transformation exceeds the range within which
a proposed invariant applies, or that the criteria
of preservation have not been adequately specified.

None of these possibilities establishes the absence
of identity.

The distinction between an invariant not being
identified and an invariant not existing must also
be maintained. Failure to establish a preservation
claim is not equivalent to establishing that no
feature could be preserved under any relevant
transformation.

Even where a particular determination is shown
to change, the ontological question of identity
remains open. Changes in composition, structure,
relations, or other determinations may be relevant
to that question, but their significance cannot
be settled merely by recording the absence of
preservation in one respect.

Whether an existent retains its identity through
a particular transformation, and under what
conditions this is possible, requires an independent
ontological analysis.

Accordingly, neither the identification nor the
non-identification of an invariant can be treated
as a conclusive determination of identity.

\subsection{Invariance as a Dimension of Ontological Characterization}

Rejecting the reduction of identity to invariance
does not render invariance ontologically insignificant.

The preservation of a determination or relation
may contribute to the characterization of an
existent. It may disclose which features remain
unchanged under a specified transformation and
which features undergo alteration. Such distinctions
can be relevant to the analysis of constitution,
composition, relational organization, and persistence.

For example, the preservation of a structural
relation may establish that a particular organizational
feature remains unchanged despite modifications
in the components participating in that relation.

The preservation of a quantitative determination
may establish that a specified magnitude remains
constant within a particular domain.

The preservation of a relation among components
may establish that a particular relational arrangement
continues to obtain despite changes in other
determinations.

In each case, the invariance claim concerns a
specified respect. Its ontological significance
depends upon the feature examined, the transformation
considered, the domain within which the claim
applies, and the criterion by which preservation
is established.

These findings may contribute to a more extensive
account of an existent. They may constrain certain
interpretations of its transformation or clarify
the respects in which its configuration is maintained
or altered.

However, characterization must not be confused
with ontological constitution. To characterize
an existent through one or more preserved features
is not necessarily to establish what constitutes
its identity.

The distinction is especially important when
invariance is used to discuss persistence. A
determination may persist, a relation may remain
preserved, or a configuration may exhibit continuity
under specified conditions. None of these descriptions
is automatically equivalent to the persistence
or continuity of the existent itself.

The relation between these concepts requires
separate examination. Invariance may contribute
to their analysis, but it cannot substitute for
their ontological clarification.

\subsection{Multiple Invariants and the Open Question of Multiple Identities}

A single existent may exhibit different forms
of invariance under different transformations
or within different domains.

Some determinations may remain preserved while
others change. A structural relation may remain
invariant under one transformation but not under
another. A relational configuration may be
preserved within a restricted domain while
undergoing alteration under a broader class
of transformations.

This plurality of invariants raises questions
concerning the relation between different respects
of preservation and the identity of the existent.

It does not, however, establish a plurality of
identities.

There is no justified inference from the existence
of several invariants to the existence of several
identities. Each invariant concerns a specified
feature under specified conditions. The number
of such features does not determine the number
of identities to which they may be related.

Nor can the totality of invariants automatically
be treated as a single identity. Their aggregation
would still leave unresolved why the selected
features should constitute identity and how their
different domains, transformations, and respects
should be related.

The possible plurality of identities is a separate
ontological problem. It cannot be settled by
counting preserved features, identifying different
forms of invariance, or assigning an identity
to every level of constitution.

Whether an existent may possess multiple identities,
whether identity is to be considered at different
ontological levels, and how any such identities
might be related remain open questions requiring
independent argument.

The analysis of invariance may help formulate
these questions more precisely. It does not
provide their answer.

\subsection{The Analytical Independence of Identity and Invariance}

The distinction developed in this section can
be summarized through the questions addressed
by each concept.

Identity concerns the ontological identity of
an existent. It asks how the existent is to be
understood in its identity, without presupposing
that this question can be answered through the
preservation of selected features.

Invariance concerns whether a specified determination,
relation, or configuration remains preserved under
a specified transformation and within a specified
domain.

The two questions are analytically independent
in the sense that neither can be answered merely
by substituting the other for it.

This independence does not imply that identity
and invariance are unrelated. An invariance claim
may contribute to the characterization of an
existent, while an account of identity may constrain
the interpretation of preservation across
configurations.

Their relation, however, requires an argument
that respects the distinct subject matter of
each inquiry.

This distinction also limits the direct transfer
of formal definitions of invariance into ontology.
A formal framework may define identity through
equivalence under a specified class of allowable
transformations. Such a definition can be valid
within the framework for which it is constructed.

Its formal validity does not, by itself, establish
that the same criterion determines the ontological
identity of an existent.

An ontological application would require independent
justification of the relevant transformations,
the domain of application, the respects of
preservation, and the reason why the resulting
equivalence relation should be interpreted as
identity rather than as a more limited form
of correspondence or sameness.

Without that justification, a formal invariance
criterion remains a framework-specific construction
and must not be elevated into a general ontological
principle.

The distinction between identity and invariance
is therefore not a denial of their possible
relation. It is a requirement that this relation
be established without reducing one problem
to the other.

\subsection{Conclusion: Invariance Without Reduction to Identity}

Identity and invariance address distinct ontological
questions.

Identity concerns the existent in its ontological
identity. Invariance concerns the possible preservation
of a specified determination, relation, or configuration
under a specified transformation and within
a specified domain.

The preservation of an invariant does not, by
itself, establish numerical identity. The absence
of an identified invariant does not establish
the absence of identity. Multiple invariants
do not entail multiple identities, and their
aggregation does not automatically constitute
a single identity.

Nevertheless, invariance remains relevant to
ontological analysis. It can characterize particular
respects of preservation, clarify the relation
between transformation and alteration, and
contribute to inquiries concerning constitution,
composition, and persistence.

Its contribution must remain proportionate to
what an invariance claim actually establishes.
A claim concerning the preservation of a feature
must not be extended into a conclusion about
the identity of the existent without an additional
argument.

The distinction established in this section
therefore preserves the ontological independence
of the problem of identity while allowing
invariance to retain its proper analytical
significance.

The next section examines the conditions under
which a determination or relation may properly
be described as invariant, including the role
of transformation, domain, respect, and the
criterion of preservation.

% =============================================================================
% END OF SECTION 02 OF CHAPTER 02
% =============================================================================